\section*{PREFACE}
\thispagestyle{empty}
% Bối cảnh và tầm quan trọng của lĩnh vực: Giới thiệu lĩnh vực nghiên cứu  Cho thấy vì sao chủ đề này quan trọng

% Nêu vấn đề: Chỉ ra khoảng trống hoặc khó khăn  -->  Trả lời câu hỏi: "Tại sao cần làm đề tài này?"

% Giải pháp và đóng góp: Đề tài làm gì Phương pháp gì Đóng góp gì

% Lời cảm ơn


As fifth-generation (5G) mobile networks continue to expand, mobility management has become an important aspect of maintaining service continuity and network performance. Among mobility procedures, handover plays a key role because user equipment may need to move between serving cells while radio conditions and network conditions change over time. This topic is important not only from the viewpoint of radio coverage, but also from the viewpoint of signaling procedures and timing behavior inside the network.

However, collecting detailed 5G mobility data in real environments is difficult because measurement campaigns require suitable equipment, stable test conditions, and repeated experiments. In addition, handover behavior cannot be fully understood from radio measurements alone, since a change in the best serving cell does not necessarily mean that the protocol-level handover procedure has been completed. These difficulties motivated the development of a controlled study environment in which radio measurements and handover-delay behavior can be inspected more systematically.

For that reason, this graduation thesis was developed to support the study of 5G mobility through simulation-based radio-data generation and log-based handover-delay analysis. The work focuses on generating structured radio-measurement data and examining handover delay from timer-event logs under different timer configurations and network-condition settings.

Through this work, radio-level serving-cell decisions and protocol-level handover completion can be considered separately. This separation provides a clearer basis for observing handover behavior and for evaluating how timer configurations and network conditions may influence completion time, timeout cases, and failure behavior.

I would like to express my deepest gratitude to \textbf{Dr. Phung Thi Kieu Ha} for her guidance, feedback, and support throughout this graduation thesis. I also thank the lecturers of the School of Electrical and Electronic Engineering, Hanoi University of Science and Technology, for providing the technical and academic foundation necessary for this work. Due to limitations in time, experimental resources, and experience, some limitations may remain. I sincerely welcome constructive comments and suggestions that can improve both the thesis and the research platform.

\clearpage
